\documentclass[12pt,a4paper]{article}

\usepackage[spanish]{babel}
\usepackage[margin=2.5cm]{geometry}
\usepackage{fontspec}
\usepackage{microtype}
\usepackage{setspace}
\usepackage{graphicx}
\usepackage{array}
\usepackage{booktabs}
\usepackage{longtable}
\usepackage{float}
\usepackage{amsmath}
\usepackage{siunitx}
\usepackage{titlesec}
\usepackage{enumitem}
\usepackage{fancyhdr}
\usepackage{hyperref}
\usepackage{xcolor}
\usepackage{caption}
\usepackage{tocloft}
\usepackage{tikz}
\usepackage[most]{tcolorbox}

\setmainfont{Times New Roman}
\setsansfont{Arial}
\setmonofont{Latin Modern Mono}

\hypersetup{
  colorlinks=true,
  linkcolor=black,
  urlcolor=black,
  pdftitle={Formato de Informe - UNI FIM},
  pdfauthor={Universidad Nacional de Ingenier\'ia},
  pdfsubject={Plantilla de informe acad\'emico de ingenier\'ia},
  pdfcreator={Codex}
}

% Paleta institucional sobria
\definecolor{UNIPrimary}{HTML}{711610}
\definecolor{UNIPrimaryDark}{HTML}{4E0F0B}
\definecolor{UNIAccent}{HTML}{B79D63}
\definecolor{UNIInk}{HTML}{1F1A17}
\definecolor{UNILine}{HTML}{D6CCC0}
\definecolor{UNILight}{HTML}{FAF8F4}

% Variables de portada
\newcommand{\reporttitle}{Informe de Laboratorio}
\newcommand{\reportsubtitle}{Ensayo Experimental de Ingenier\'ia Mec\'anica}
\newcommand{\reportcourse}{Nombre del curso}
\newcommand{\reportauthor}{Nombre del estudiante}
\newcommand{\reportcode}{C\'odigo universitario}
\newcommand{\reportteacher}{Nombre del docente}
\newcommand{\reportgroup}{Secci\'on o grupo}
\newcommand{\reportlab}{Laboratorio o asignatura}
\newcommand{\reportdate}{06/04/2026}
\newcommand{\reportterm}{2026-I}
\newcommand{\reportcity}{Lima - Per\'u}
\newcommand{\unishieldpath}{escudo_uni_portada_crisp.png}

\setstretch{1.18}
\setlength{\parindent}{1.1em}
\setlength{\parskip}{0.3em}
\setlength{\headheight}{24.1pt}

% Encabezado y pie
\pagestyle{fancy}
\fancyhf{}
\fancyhead[L]{\small\textcolor{UNIPrimaryDark}{Universidad Nacional de Ingenier\'ia}}
\fancyhead[R]{\small\textcolor{UNIPrimaryDark}{Facultad de Ingenier\'ia Mec\'anica}}
\fancyfoot[L]{\small\textcolor{UNIPrimaryDark}{\reporttitle}}
\fancyfoot[R]{\small\textcolor{UNIPrimaryDark}{\thepage}}
\renewcommand{\headrulewidth}{0.5pt}
\renewcommand{\headrule}{\hbox to\headwidth{\color{UNILine}\leaders\hrule height \headrulewidth\hfill}}
\renewcommand{\footrulewidth}{0pt}

% Jerarquia tipografica
\titleformat{\section}
  {\Large\bfseries\color{UNIPrimaryDark}}
  {\thesection}
  {0.65em}
  {}

\titleformat{\subsection}
  {\large\bfseries\color{UNIPrimaryDark}}
  {\thesubsection}
  {0.65em}
  {}

\titlespacing*{\section}{0pt}{1.2em}{0.55em}
\titlespacing*{\subsection}{0pt}{0.95em}{0.35em}

% TOC
\renewcommand{\cfttoctitlefont}{\Huge\bfseries\color{UNIPrimaryDark}}
\renewcommand{\cftsecfont}{\color{UNIInk}}
\renewcommand{\cftsecpagefont}{\color{UNIInk}}
\setlength{\cftbeforesecskip}{0.25em}

% Figuras y tablas
\captionsetup{
  font=small,
  labelfont={bf,color=UNIPrimaryDark},
  textfont=small
}
\sisetup{
  locale = DE,
  output-decimal-marker = {,},
  per-mode = symbol
}

\newtcolorbox{infobox}[1][]{
  enhanced,
  colback=UNILight,
  colframe=UNIAccent,
  boxrule=0.6pt,
  arc=1.2mm,
  left=4.5mm,
  right=4.5mm,
  top=2.2mm,
  bottom=2.2mm,
  #1
}

\newcommand{\EscudoUNI}{%
  \IfFileExists{\unishieldpath}{%
    \includegraphics[width=5.1cm]{\unishieldpath}%
  }{%
    \begin{tikzpicture}
      \draw[UNIPrimaryDark,line width=0.9pt] (0,0) circle (1.5cm);
      \draw[UNIAccent,line width=0.7pt] (0,0) circle (1.3cm);
      \node[align=center,text width=2.3cm,font=\sffamily\bfseries\small,text=UNIPrimaryDark] at (0,0.1) {ESCUDO\\UNI};
    \end{tikzpicture}%
  }%
}

\newcommand{\PortadaUNI}{%
\begin{titlepage}
  \thispagestyle{empty}
  \color{UNIInk}

  \begin{center}
    {\bfseries\fontsize{18pt}{21pt}\selectfont UNIVERSIDAD NACIONAL DE\\INGENIER\'IA}\\[0.28cm]
    {\fontsize{12.5pt}{14pt}\selectfont FACULTAD DE INGENIER\'IA MEC\'ANICA}
  \end{center}

  \vspace{1.15cm}
  \begin{center}
    \EscudoUNI
  \end{center}

  \vspace{0.85cm}
  \begin{center}
    {\bfseries\fontsize{12.8pt}{14.8pt}\selectfont \MakeUppercase{\reportlab}}\\[0.28cm]
    {\bfseries\fontsize{14pt}{16pt}\selectfont \MakeUppercase{\reporttitle}: \MakeUppercase{\reportsubtitle}}
  \end{center}

  \vspace{1.05cm}
  \begin{center}
    \renewcommand{\arraystretch}{1.72}
    \begin{tabular}{>{\bfseries}p{3.4cm} p{8.1cm}}
      CURSO: & \reportcourse \\
      SECCI\'ON: & \reportgroup \\
      DOCENTE: & \reportteacher \\
      ESTUDIANTE: & \reportauthor \\
      C\'ODIGO: & \reportcode \\
      FECHA: & \reportdate \\
    \end{tabular}
  \end{center}

  \vfill
  \begin{center}
    {\fontsize{23pt}{26pt}\selectfont \reportterm}\\[0.38cm]
    {\small \reportcity}
  \end{center}
\end{titlepage}
}

\begin{document}

\PortadaUNI

\tableofcontents
\thispagestyle{empty}
\clearpage

\section{Introducci\'on}
Describe el contexto del ensayo, el problema que motiva el informe y su relevancia dentro del curso. Cierra esta secci\'on con una descripci\'on breve de la estructura del documento.

\section{Objetivos}
\begin{itemize}[leftmargin=1.5em]
  \item Objetivo general del ensayo o pr\'actica.
  \item Objetivo espec\'ifico 1.
  \item Objetivo espec\'ifico 2.
  \item Objetivo espec\'ifico 3.
\end{itemize}

\section{Fundamento Te\'orico}
Incluye conceptos, principios f\'isicos, ecuaciones y criterios t\'ecnicos necesarios para interpretar el ensayo.

\subsection{Principios relevantes}
Explica aqu\'i las leyes, modelos o definiciones que sustentan el trabajo.

\subsection{Ecuaciones base}
\[
  \eta = \frac{P_{\text{\'util}}}{P_{\text{entrada}}}\times 100
\]
\[
  Q = m\,c\,\Delta T
\]

\section{Equipos}
Enumera y describe los equipos, instrumentos y materiales utilizados.

\begin{longtable}{p{4.5cm} p{4cm} p{5cm}}
\toprule
\textbf{Equipo o instrumento} & \textbf{Modelo / rango} & \textbf{Observaci\'on} \\
\midrule
Banco de ensayo & --- & Completar \\
Sensor de temperatura & --- & Completar \\
Tac\'ometro / mult\'imetro & --- & Completar \\
\bottomrule
\end{longtable}

\section{Procedimiento del Ensayo}
Presenta los pasos ejecutados en orden cronol\'ogico y con detalle suficiente para permitir la repetici\'on del ensayo.

\begin{enumerate}[leftmargin=1.5em]
  \item Preparaci\'on del banco o sistema experimental.
  \item Verificaci\'on de condiciones iniciales y calibraci\'on de instrumentos.
  \item Toma de datos bajo diferentes condiciones de operaci\'on.
  \item Registro de incidencias durante la prueba.
\end{enumerate}

\section{Datos del Laboratorio}
Registra las mediciones obtenidas directamente del experimento.

\begin{table}[H]
  \centering
  \renewcommand{\arraystretch}{1.25}
  \begin{tabular}{>{\bfseries}c c c c c}
    \toprule
    Ensayo & Variable 1 & Variable 2 & Variable 3 & Observaci\'on \\
    \midrule
    1 & --- & --- & --- & --- \\
    2 & --- & --- & --- & --- \\
    3 & --- & --- & --- & --- \\
    \bottomrule
  \end{tabular}
  \caption{Datos experimentales iniciales.}
\end{table}

\section{C\'alculos y Resultados}
Desarrolla los c\'alculos a partir de los datos medidos. Conviene separar claramente desarrollo y resultados.

\subsection{C\'alculos}
Muestra una sustituci\'on ordenada de valores, unidades y resultados intermedios.

\begin{infobox}
Ejemplo de presentaci\'on:
\[
  P = T\,\omega
\]
Si $T = \SI{12}{N.m}$ y $\omega = \SI{18}{rad/s}$, entonces:
\[
  P = 12\times 18 = \SI{216}{W}
\]
\end{infobox}

\subsection{Resultados}
Resume los resultados m\'as importantes y acompa\~nalos con tablas, figuras o comparaciones te\'orico-experimentales.

\begin{table}[H]
  \centering
  \renewcommand{\arraystretch}{1.25}
  \begin{tabular}{>{\bfseries}p{5cm} p{4cm} p{4cm}}
    \toprule
    Magnitud & Valor obtenido & Unidad \\
    \midrule
    Eficiencia & --- & \% \\
    Potencia & --- & W \\
    Error relativo & --- & \% \\
    \bottomrule
  \end{tabular}
  \caption{Resumen de resultados finales.}
\end{table}

\section{Observaciones}
Registra condiciones reales del ensayo, limitaciones, desviaciones del procedimiento y cualquier aspecto que influya en la interpretaci\'on de los datos.

\section{Conclusiones}
\begin{itemize}[leftmargin=1.5em]
  \item Conclusi\'on principal alineada con el objetivo general.
  \item Interpretaci\'on t\'ecnica del resultado m\'as relevante.
  \item Recomendaci\'on para mejorar el procedimiento en pruebas futuras.
\end{itemize}

\section{Bibliograf\'ia (APA 7.\textordfeminine~edici\'on)}
Usa formato APA para libros, art\'iculos, normas y fuentes web. Referencias de ejemplo:

\begin{thebibliography}{9}
  \bibitem{manual_lab}
  Autor, A. A. (2022). \textit{Manual del laboratorio o gu\'ia de pr\'actica}. Editorial o Instituci\'on.

  \bibitem{mecanica}
  Autor, B. B., \& Autor, C. C. (2021). \textit{Texto de ingenier\'ia mec\'anica} (3.\textordfeminine~ed.). Editorial.

  \bibitem{norma}
  American Psychological Association. (2020). \textit{Publication manual of the American Psychological Association} (7th ed.).
\end{thebibliography}

\section{Anexos}
Incluye fotograf\'ias, hojas de datos, tablas extendidas, capturas del banco de pruebas, planos o c\'alculos complementarios.

\begin{infobox}[colback=white,colframe=UNILine]
Sugerencias de uso:
\begin{itemize}[leftmargin=1.5em]
  \item Reemplaza los campos de la portada al inicio del archivo.
  \item Inserta figuras con \texttt{\textbackslash includegraphics}.
  \item Mant\'en unidades normalizadas (SI) y consistencia decimal en tablas.
\end{itemize}
\end{infobox}

\end{document}
